% Generated draft: requires mathematical and source review.
\documentclass[11pt]{article}
\usepackage[margin=25mm]{geometry}
\usepackage{fontspec}
\setmainfont{Latin Modern Roman}
\usepackage{amsmath,amssymb,mathtools,unicode-math}
\setmathfont{Latin Modern Math}
\usepackage{enumitem,needspace,xurl,hyperref}
\setlength{\parindent}{0pt}
\setlength{\parskip}{5pt}
\setlength{\emergencystretch}{3em}
\newcommand{\N}{\mathbb{N}}
\newcommand{\Z}{\mathbb{Z}}
\newcommand{\Q}{\mathbb{Q}}
\newcommand{\R}{\mathbb{R}}
\newcommand{\C}{\mathbb{C}}
\newcommand{\F}{\mathbb{F}}
\newcommand{\A}{\mathbb{A}}
\newcommand{\PP}{\mathbb{P}}
\newcommand{\E}{\mathbb{E}}
\newcommand{\eps}{\varepsilon}
\newcommand{\dd}{\,\mathrm{d}}
\begin{document}
\section*{UnsolvedMath ID: 800510}
\textbf{Source problem code:} catalogue entry
\subsection{Definitions and mathematical statement}
The supplied record asks for an extension of Paige's analysis of the scalar Lanczos algorithm, cited there as [paige\_76, paige\_80], to the block Lanczos algorithm. A block formulation processes several vectors together; orthogonalization within a block is the procedure used to produce mutually orthonormal vectors spanning that block's retained subspace. The scalar field, inner product, treatment of dependent vectors, and precise block recurrence must be specified in a complete formulation.

\textbf{Requested research statement.} Develop a block Lanczos analogue of the analysis in [paige\_76, paige\_80]. The analysis must account explicitly for the accuracy of orthogonalization within each block. As required by the source, either state assumptions on the accuracy of that orthogonalization procedure and develop the analysis under those assumptions, or specify and analyze a standard algorithm that satisfies an appropriate, explicitly stated interpretation of the source's description ``good.''

\textbf{Unresolved qualifications.} The supplied record does not reproduce Paige's results or identify the particular conclusions to be extended. It also does not specify the matrix class, scalar field, arithmetic or error model, block sizes, initialization, recurrence, interblock orthogonalization policy, stopping conditions, or treatment of rank deficiency and breakdown. In particular, no matrix symmetry or self-adjointness hypothesis is explicitly supplied. The source gives no quantitative definition of accuracy or of a ``good'' orthogonalization algorithm. These choices and the intended analogue of Paige's conclusions must be fixed before the request can be expressed as a fully quantified theorem or numerical error bound. No particular choices are imposed in this draft.
\subsection{Short English statement}
Extend Paige's cited Lanczos analysis to block Lanczos, explicitly accounting for orthogonalization within blocks. Either formulate accuracy assumptions for that procedure or analyze a specified standard algorithm. The supplied source leaves the precise conclusions, hypotheses, and numerical model unresolved.
\subsection{Sources}
\begin{itemize}
\item Local UnsolvedMath registry, permanent ID \texttt{800510}, source problem code \texttt{catalogue entry}, title ``Linear Systems and Eigenvalue Problems --- Problem 19.'' The registry attributes the problem to Amsel, Baumann et al., with proposed year 2026, and labels it \emph{open}; this is supplied status metadata, not an independently verified mathematical or literature assessment.
\item Amsel, Baumann et al., \emph{Linear Systems and Eigenvalue Problems: Open Questions from a Simons Workshop} (2026), Problem 19. Supplied source URL: \url{https://arxiv.org/abs/2602.05394}. The record reports access on 2026-10-16 and extraction by \texttt{arxiv\_latex}; those are supplied provenance statements, not a claim of access or verification for this draft.
\item Paige references [paige\_76] and [paige\_80], as cited in the supplied statement. Full bibliographic entries and the relevant results are not supplied, so these references cannot be identified precisely from this record alone.
\end{itemize}
\end{document}
